% ============================================================
%  chapters/chapter3.tex
%  Implementation and Validation
%  COMP3931 Individual Project  Adhiraj Kumar  2025/26
% ============================================================

\chapter{Implementation and Validation}
\label{ch:implementation}

Chapter~2 discussed the design choices. Each detector was selected for a reason, every parameter was justified from our principles, and project risks were anticipated. Naturally the next question is whether the system built from those choices actually works. This chapter answers that through a overview of the five components, detailed detector implementations, a description of the controller, and a layered validation programme that establishes correctness before any real data experiments are run.

% -----------------------------------------------------------
\section{System Overview}
\label{sec:sys_overview}
% -----------------------------------------------------------

The framework comprises five loosely coupled components, each with a single responsibility. Loose coupling means each component can be tested, modified, and replaced independently. The full pipeline is shown in Figure~\ref{fig:architecture} (Appendix~B).

\begin{enumerate}
  \item \textbf{Data layer} (\texttt{src/data\_loader.py}): downloads OHLCV and FRED
    macro series, constructs the 33 feature matrix, and encodes the binary target.
  \item \textbf{Model layer} (\texttt{src/model/}): maintains static and adaptive
    instances under a common interface for RF, GBM, and LR, each wrapped in a
    \texttt{RobustScaler} scikit learn \texttt{Pipeline}.
  \item \textbf{Detector stack} (\texttt{src/drift\_detectors/}): five detector
    classes exposing a common \texttt{detect(reference, current)} interface. PH is
    the primary performance monitor. KS, PSI, JS, and the prediction drift monitor
    are supplementary probes testing the lead time hypothesis from Chapter~2,
    Section~\ref{sec:detector_rationale}.
  \item \textbf{Controller} (\texttt{src/controller/}): aggregates scores into a
    composite drift index, calibrates asset adaptive thresholds, and dispatches one
    of four retraining actions.
  \item \textbf{Evaluation layer} (\texttt{src/evaluation/}): records per run
    outcomes and contains the 23 experiment scripts used in Chapter~4.
\end{enumerate}

Detectors produce raw numerical scores. The controller handles thresholding, weighting, and dispatch, keeping each component independently testable.

% -----------------------------------------------------------
\section{Drift Detector Implementations}
\label{sec:detectors}
% -----------------------------------------------------------

Each detector is implemented as an original Python class from the mathematical definitions in the cited papers. No detection logic was copied from external libraries. The PSI implementation reproduces the worked example in Siddiqi~\cite{siddiqi2006} to four decimal places as a correctness check.

\subsection{Kolmogorov Smirnov Test}
\label{subsec:ks}

The KS statistic measures the largest vertical gap between two empirical cumulative distribution functions (CDFs). An empirical CDF is the step function that rises by $1/n$ at each observed data point. The gap between two CDFs at any given point tells you how differently the two samples are distributed up to that value.

To make this concrete: imagine plotting the sorted feature values from a reference window and from the current window as two staircase curves. The KS statistic is the maximum height difference between those two staircases anywhere along the x axis. A large gap means the current distribution has shifted substantially from the reference:

\begin{equation}
  D_{n,m} = \sup_{x} \left| F_n(x) - G_m(x) \right|
  \label{eq:ks}
\end{equation}

where $F_n(x)$ is the reference window CDF and $G_m(x)$ is the current window CDF. $D_{n,m} \in [0,1]$, non parametric~\cite{kolmogorov1933,smirnov1948,engle1982}. The implementation uses a merge walk over two sorted arrays in $\mathcal{O}(n \log n)$. $p$ values are from \texttt{scipy.stats.ks\_2samp}. KS is applied per feature, and the maximum across all features is the aggregate. The ablation in Chapter~4 (Table~B.1, Appendix~B) confirms KS is redundant when PSI is active~\cite{rabanser2019}.

\subsection{Population Stability Index}
\label{subsec:psi}

PSI measures how much a population has shifted relative to a reference by comparing the proportion of observations falling in each quantile bin. The intuition is straightforward: if 20\% of reference data fell in a certain range and now only 5\% does, the population has moved away from that range. PSI captures this movement across all bins simultaneously and sums it into a single score. Because the bins are defined by the reference distribution's own deciles, each bin holds roughly equal mass in the reference window, making PSI sensitive to shifts anywhere in the distribution and not just in the tails:

\begin{equation}
  \text{PSI} = \sum_{b=1}^{B} \left( p_b^{\text{current}} - p_b^{\text{ref}} \right)
    \ln \frac{p_b^{\text{current}}}{p_b^{\text{ref}}}
  \label{eq:psi}
\end{equation}

where $B = 10$ decile bins are defined by the reference window quantiles. $p_b^{\text{ref}}$ and $p_b^{\text{current}}$ are the corresponding bin fractions. Laplace smoothing ($\epsilon = 10^{-6}$) prevents division by zero. PSI is applied per feature. The maximum is the aggregate.

\subsection{Jensen Shannon Divergence}
\label{subsec:js}

JS divergence is applied to the model's output probability distributions, monitoring prediction drift. To understand why JS rather than the more familiar KL divergence: KL divergence measures how much one distribution differs from another, but it is asymmetric ($D_{\text{KL}}(P \| Q) \neq D_{\text{KL}}(Q \| P)$) and can be infinite if one distribution assigns zero probability where the other does not. JS solves both problems by comparing each distribution not directly to the other, but to their mixture. The result is always finite, always symmetric, and bounded in $[0, \ln 2]$, making it safe to normalise for the composite index:

\begin{equation}
  \text{JS}(P \| Q) = \frac{1}{2} D_{\text{KL}}\!\left(P \,\Big\|\, \frac{P+Q}{2}\right)
                    + \frac{1}{2} D_{\text{KL}}\!\left(Q \,\Big\|\, \frac{P+Q}{2}\right)
  \label{eq:js}
\end{equation}

where $P$ is the reference window output distribution, $Q$ is the current window output distribution, and $D_{\text{KL}}(A \| B) = \sum_i A_i \ln(A_i / B_i)$ is the Kullback Leibler divergence~\cite{lin1991,shannon1948,kullback1951}. Both distributions are discretised into 20 equal width bins. Laplace smoothing prevents log of zero at bin boundaries.

\subsection{Page Hinkley Sequential Change Point Detector}
\label{subsec:ph}

PH~\cite{page1954,mouss2004} detects a persistent upward shift in the mean of a stream. Applied here, it monitors rolling log loss. Log loss penalises confident wrong predictions heavily: if the model assigns a 95\% probability to an up day and the day turns out to be a down day, the penalty is severe. A sustained rise in log loss therefore signals that the model's calibration is degrading and that its confident predictions are increasingly pointing the wrong way.

PH works by accumulating the deviation of each new observation from the running mean. Rather than reacting to any single large value, which could just be a fat tailed return day, it requires the deviations to compound over time before raising an alarm. This persistence requirement is what makes it reliable in financial data, where individual extreme days are common but sustained deterioration is the signal worth catching. State updates at each observation $x_t$:

\begin{align}
  \bar{x}_t  &= \frac{1}{t} \sum_{i=1}^{t} x_i  \label{eq:ph_mean} \\
  U_t        &= U_{t-1} + (x_t - \bar{x}_t - \delta) \label{eq:ph_sum} \\
  m_t        &= \min(U_1, \ldots, U_t) \label{eq:ph_min} \\
  \text{PH}_t &= U_t - m_t \label{eq:ph_stat}
\end{align}

where $\delta > 0$ is an asset adaptive sensitivity parameter that prevents small fluctuations from accumulating in $U_t$. The alarm fires when $\text{PH}_t$ exceeds the per asset threshold $\lambda$. Each update requires $\mathcal{O}(1)$ time and space, maintaining five scalar state variables. PH is the only detector suitable for genuine streaming at tick resolution.

\subsection{Prediction Drift Monitor}
\label{subsec:pred_drift}

The prediction drift monitor applies JS divergence to the model's output probability vectors rather than input features. It catches a specific failure mode that neither KS nor PSI would detect: the case where features more extreme than any in the training set push the model toward clustered near zero or near one outputs, without a corresponding shift in the marginal feature distributions.

When this happens, the model is overconfident. It assigns probabilities close to 0 or 1 on days where it should be uncertain. But the feature distributions still look stable, because the extreme values are rare enough not to shift the population statistics. JS on the output distribution catches this directly: a shift from a well spread output distribution toward a bimodal one near the extremes produces a large divergence from the reference window output profile, even when KS and PSI are quiet. This makes the prediction drift monitor an orthogonal signal to the feature distribution stack, not a redundant one.

% -----------------------------------------------------------
\section{Drift Controller and Calibration}
\label{sec:controller}
% -----------------------------------------------------------

\subsection{Composite Drift Index}
\label{subsec:composite}

The drift controller (\texttt{src/controller/drift\_controller.py}) combines detector outputs into $\mathcal{D}_t$ as defined in Chapter~2 (Equation~\ref{eq:composite_ch2}). $s_{\text{feature}}$ is the maximum normalised KS/PSI score across features. $s_{\text{prediction}}$ is the normalised JS score from the prediction drift monitor. $s_{\text{performance}}$ is the normalised PH statistic. All scores are normalised to $[0,1]$ before weighting. Both $\mathcal{D}_t$ and $s_{\text{performance}}$ are EMA smoothed at $\alpha{=}0.10$ before the threshold comparator.

\subsection{Asset Adaptive Calibration}
\label{subsec:calibration}

The calibration module (\texttt{src/controller/calibration.py}) learns asset specific thresholds from each instrument's reference window at the start of each evaluation run. The motivation is simple: what counts as an anomalous drift score for one instrument can be completely routine for another. Bitcoin exhibits daily moves of several hundred basis points as a matter of course. The same move in a US Treasury bond fund is a significant rate shock. A single global threshold applied to both would either fire constantly on Bitcoin (too sensitive) or miss shocks in fixed income (too conservative). Asset adaptive calibration solves this by anchoring each threshold to the instrument's own historical drift distribution. A 90th percentile event in Bitcoin therefore corresponds to an unusual day for Bitcoin, not just a normal volatile day expressed in the same numerical units as a quiet day for TLT. The mapping is:

\begin{equation}
  \tau_k = \text{Percentile}_{p_k}\!\left(\tilde{\mathcal{D}}_{1}, \ldots,
    \tilde{\mathcal{D}}_{T_{\text{ref}}}\right), \quad
  (p_1, p_2, p_3, p_4) = (60, 75, 90, 95)
  \label{eq:calibration}
\end{equation}

\subsection{Tiered Action Dispatch}
\label{subsec:tiered_actions}

The controller maps the smoothed composite drift index to one of four tiers, each dispatching a proportionately stronger response:

\begin{table}[htbp]
  \centering
  \caption{Controller action tiers. Thresholds $\tau_1$ to $\tau_4$ are set per asset
    by calibration. Default values apply when the reference window is too short.}
  \label{tab:controller_tiers}
  \begin{tabular}{llll}
    \hline
    \textbf{Tier} & \textbf{Condition} & \textbf{Default threshold}
      & \textbf{Action dispatched} \\
    \hline
    None     & $\tilde{\mathcal{D}}_t < \tau_1$ & $< 0.6$  & No action \\
    Low      & $\tau_1 \leq \tilde{\mathcal{D}}_t < \tau_2$ & $[0.6, 1.0)$ & EMA position update only \\
    Moderate & $\tau_2 \leq \tilde{\mathcal{D}}_t < \tau_3$ & $[1.0, 1.5)$ & SGD weighted update \\
    High     & $\tau_3 \leq \tilde{\mathcal{D}}_t < \tau_4$ & $[1.5, 2.0)$ & Sliding window retrain \\
    Severe   & $\tilde{\mathcal{D}}_t \geq \tau_4$          & $\geq 2.0$   & Ensemble refresh \\
    \hline
  \end{tabular}
\end{table}

% -----------------------------------------------------------
\section{Adaptation Strategies}
\label{sec:adaptation}
% -----------------------------------------------------------

Three strategies are dispatched at Moderate, High, and Severe tiers (\texttt{src/controller/adaptation.py}). The Low tier adjusts only the position gate.

\paragraph{Weighted ensemble blend (Moderate tier).}
Rather than discarding prior learning, the model is nudged toward recent data:

\begin{equation}
  \theta_{\text{new}} = (1 - \gamma)\,\theta_{\text{old}} + \gamma\,\theta_{\text{recent}},
  \quad \gamma = 0.15
  \label{eq:sgd_blend}
\end{equation}

The 85\%/15\% split reflects the stability plasticity balance of Herbster and Warmuth~\cite{herbster1998}. Ablation support is in Chapter~4, Table~\ref{tab:retrain_summary}.

\paragraph{Sliding window retrain (High tier).}
The model is fully retrained on the most recent 100 trading days. The \texttt{min\_retrain\_rows = 400} constraint is enforced. If fewer than 400 observations are available, the retrain is deferred.

\paragraph{Ensemble refresh (Severe tier).}
The model is reinitialised from scratch and retrained on the most recent available window subject to the 400 row minimum. If the refreshed model's accuracy falls more than 8 percentage points below the static model, the controller reverts to the pre retrain snapshot.

% -----------------------------------------------------------
\section{Position Gated Blend}
\label{sec:position_gate}
% -----------------------------------------------------------

The position gate blends adaptive and static outputs continuously rather than switching between them, avoiding calibration errors from a hard switch:

\begin{equation}
  \hat{p}_t = w_{\text{adapt}} \cdot \hat{p}_{\text{adaptive},t}
            + (1 - w_{\text{adapt}}) \cdot \hat{p}_{\text{static},t}
  \label{eq:blend}
\end{equation}

where $w_{\text{adapt}} \in [0.25, 0.90]$ is computed from the gate formula (Chapter~2, Equation~\ref{eq:gate_ch2}). The clip bounds ensure the adaptive model always contributes at least 25\% and the static model retains at least 10\%.

% -----------------------------------------------------------
\section{Testing and Validation}
\label{sec:testing}
% -----------------------------------------------------------

Three validation layers ensure any result in Chapter~4 reflects framework behaviour, not an implementation artefact. 


\subsection{Unit Tests}
\label{subsec:unit_tests}

The test suite contains 121 tests across 5 modules, all passing. The five detector classes, the controller, and the calibration module achieve 100\% line coverage. Four property categories are highlighted.

\begin{itemize}
  \item \textbf{Look ahead test.} The feature value on day $t$ is unchanged when
    price data on days $t{+}1$ through $t{+}5$ is perturbed, verified across all 33
    features. An inadvertent look ahead bias would produce unrealistically high
    backtested accuracy and invalidate all downstream results.

  \item \textbf{Page Hinkley alarm behaviour.} Three properties are verified: (i)
    $\text{PH}_t = 0$ at initialisation; (ii) no alarm on a stationary Gaussian
    sequence at the default threshold; (iii) an alarm fires within a bounded number
    of steps after $\Delta\mu \geq 0.5\sigma$.

  \item \textbf{Calibration monotonicity and fallback.} Thresholds satisfy
    $\tau_1 \leq \tau_2 \leq \tau_3 \leq \tau_4$. Calibration is idempotent. The
    fallback triggers correctly when the reference window is below minimum size.

  \item \textbf{Position gate clip bounds.} $w_{\text{adapt}} \in [0.25, 0.90]$
    for all inputs. Blended probability $\hat{p}_t \in [0,1]$.
\end{itemize}

\subsection{Synthetic Benchmark}
\label{subsec:synth_validation}

Without known change points, the true positive rate (TPR, the fraction of genuine drift events that trigger an alarm) and false positive rate (FPR, the fraction of drift free days that produce a spurious alarm) cannot be computed from real data. Controlled Gaussian sequences with a known mean shift $\Delta\mu$ at a fixed position provide exact ground truth. Three magnitudes ($\Delta\mu \in \{0.5\sigma, 1.0\sigma, 2.0\sigma\}$) were tested with 20 replications each (approximately 420 runs).

At $\Delta\mu \geq 1.0\sigma$, PSI+PH achieved TPR 93 to 100\% with FPR below 5\% and median alarm latency of zero observations. ADWIN~\cite{bifet2007adwin} achieved 0.2\% TPR under identical conditions. This confirms quantitatively that error stream detectors cannot catch a feature level shift before it propagates to accuracy, the same constraint that applies to DDM~\cite{gama2004ddm} and EDDM~\cite{baenagarcia2006eddm}. A supplementary Student $t$ leptokurtic benchmark, where the $t$ distribution's heavier than Gaussian tails generate occasional extreme observations at the rate typical of financial returns (Figure~\ref{fig:synth_appendix}, Appendix~B), confirms comparable TPR under heavy tailed distributions. Full results across all 16 configurations are in Chapter~4, Section~\ref{subsec:synthetic}.

\subsection{End to End Integration Test}
\label{subsec:e2e}

An end to end test on AAPL 2018 to 2022 verifies five bookkeeping properties: (i) no NaN or infinite feature values; (ii) static model parameters unchanged; (iii) retrain history correctly logged with dates, tiers, and pre and post accuracy; (iv) position gate $\hat{p}_t \in [0,1]$; (v) Sharpe annualisation uses 252 trading days, multiplying the daily Sharpe ratio by $\sqrt{252}$ to convert to an annual scale. All five pass. The test yields an accuracy delta of $+0.010$ and Sharpe delta of $+0.73$ over 29 retrains. These are reported as bookkeeping checks only. The elevated Sharpe reflects the COVID 19 dislocation in this window.

\subsection{Ablation Study Design}
\label{subsec:ablation_design}

With correctness established, two ablations are pre registered before the main experiments run.

\textbf{Detector ablation.} All $2^4{=}16$ feature detector combinations (prediction drift monitor active throughout) plus a pure static baseline, across 3,332 runs (50 instruments $\times$ 4 sub periods $\times$ 17 configurations). This tests the lead time hypothesis: if upstream KS/PSI/JS detectors fire before PH, configurations including them should outperform PH alone. The Pareto frontier, meaning the set of configurations where no alternative achieves better accuracy without also increasing runtime or vice versa, identifies the minimum sufficient configuration.

\textbf{Retrain strategy ablation.} Five strategies across 980 runs (50 instruments $\times$ 4 sub periods $\times$ 5 variants), testing whether adding an error rate trigger to distribution based detection provides any incremental benefit. Results are in Chapter~4, Section~\ref{sec:rq3}.

% -----------------------------------------------------------
\section{Monitoring Dashboard and API}
\label{sec:dashboard}
% -----------------------------------------------------------

A real time monitoring dashboard, built with a FastAPI backend and a React/TypeScript frontend, exposes drift scores, retrain history, and rolling performance via five documented REST endpoints. Examiners can inspect drift events and rolling performance interactively without re running experiment scripts. It is not used in the quantitative evaluation but demonstrates the framework is production deployable rather than purely a research artefact.

The framework is now implemented and validated. Chapter~4 presents what happens when it runs on real market data across 4,312 experimental configurations.